\section{Thesis Contributions}

This work has three main goals: research, design, and implementation. Together, they address the infrastructure sprawl problem identified in the motivation: small academic departments need a \emph{single}, unified platform that is both feature-rich enough to support modern AI/ML workflows and lightweight enough to be operated by a single administrator on consumer-grade hardware.

\subsection{Research}

The research goal is to conduct an analysis of compute resource usage at our department. This includes describing different use-cases users have for the system in place and creating typical usage scenarios for both users and administrators. User interviews were conducted with students experienced with the current solution, covering their workflows, data, experiment outputs, hardware and software needs, limiting factors, and pain points.

\todo{Forward-reference the chapter/section where interview findings are presented in detail.}

\subsection{Design}

The design goal is to design solutions that support the identified use-cases. The solution should:
\begin{itemize}
    \item be able to exist within the existing infrastructure,
    \item be easily scalable,
    \item provide additional services,
    \item support modern AI/ML practices,
    \item be maintainable long-term by a \textbf{single administrator}, thereby minimizing the total cost of ownership,
    \item \textbf{consolidate batch and interactive workloads} onto a single control plane, eliminating the need for separate HPC and cloud-native stacks,
    \item operate efficiently on \textbf{consumer-grade hardware} without relying on enterprise features such as NVIDIA MIG for GPU partitioning, InfiniBand for high-speed interconnects, or ECC memory.
\end{itemize}

\subsection{Implementation}

After providing solutions that address users' needs, the goal is to select a solution and implement it. Specifically, this thesis contributes:

\begin{itemize}
    \item A modular, Kubernetes-native platform for AI research.
    \item An intelligent scheduling layer for resource-efficient batch processing.
    \item A distributed executor for scaling research workloads.
    \item A converged, lightweight control plane, built on K3s, that unifies batch scheduling, interactive development environments, and experiment tracking with minimal operational overhead, designed to be managed by a single administrator on consumer-grade hardware.
\end{itemize}
